\documentclass[12pt, a4paper]{article}


% --- Encoding and language ---
\usepackage[utf8]{inputenc}
\usepackage[T1]{fontenc}
\usepackage[provide=*,spanish]{babel}


% --- Math ---
\usepackage{amsmath}
\usepackage{amssymb}


% --- Figures ---
\usepackage{graphicx}
\usepackage{float}
\graphicspath{{images/}}


% Gracefully show a placeholder when a figure file is missing
\newcommand{\safeincludegraphics}[1]{%
   \IfFileExists{images/#1}{%
       \includegraphics[width=\textwidth]{#1}%
   }{%
       \begin{center}%
       \fbox{\parbox{0.75\textwidth}{\centering\small\itshape%
           Figura \texttt{#1} pendiente --- ejecutar \texttt{main.ipynb} para generarla.}}%
       \end{center}%
   }%
}


% --- Tables ---
\usepackage{booktabs}


% --- Links ---
\usepackage[colorlinks=true, linkcolor=blue, citecolor=blue, urlcolor=blue]{hyperref}


% --- Page layout ---
\usepackage{geometry}
\geometry{margin=2.5cm}


% --- Header/footer ---
\usepackage{fancyhdr}
\setlength{\headheight}{14.5pt}
\addtolength{\topmargin}{-2.5pt}
\pagestyle{fancy}


\fancyhf{}
\rhead{Parcial 2}
\lhead{Thomas Gonzalez}
\cfoot{\thepage}


% ============================================================
\begin{document}


\begin{center}
   {\LARGE\bfseries Parcial 2}\\[0.8em]
   {\large Thomas Gonzalez}\\[0.3em]
   {\normalsize \today}\\[0.4em]
   {\small \href{https://github.com/thomatrixting/solar_esturcture_modeling}{github del proyecto}}
\end{center}


\vspace{1em}
\hrule
\vspace{1.5em}


% ============================================================
\section*{Contexto de los datos utilizados}


Los datos utilizados provienen del Modelo S de Jørgen Christensen-Dalsgaard et al \cite{ChristensenDalsgaard1996}.
El conjunto de datos contiene un perfil radial del Sol compuesto por 2482 capas, con las siguientes variables: radio normalizado $r/R_\odot$, velocidad del sonido $c$, densidad $\rho$, presión $P$, exponente adiabático $\Gamma_1$ y temperatura $T$.
Todas las constantes físicas empleadas en los cálculos (constante gravitacional $G$, constante de Boltzmann $k_B$, radio y masa solares $R_\odot$ y $M_\odot$, masa del protón $m_p$, constante de Stefan-Boltzmann $\sigma$ y velocidad de la luz $c$) se obtuvieron de la librería \texttt{astropy}, la cual implementa los valores recomendados por CODATA 2014 \cite{mohrCODATARecommendedValues2016}.


% ============================================================
\section*{Teoría}


\subsection*{Distribución de masa y energía potencial gravitacional}


La \textbf{ecuación de continuidad de la masa} establece que la masa encerrada en una cáscara esférica de radio $r$ y grosor $dr$ es:
\begin{equation}
   dM(r) = 4\pi r^2 \rho(r)\, dr,
\end{equation}
de modo que la masa total encerrada hasta el radio $r$ se obtiene integrando desde el centro:
\begin{equation}
   M(r) = \int_0^r 4\pi r'^2 \rho(r')\, dr'.
\end{equation}
En la práctica, esta integral se discretiza sumando los incrementos $\Delta M_i = 4\pi r_i^2\,\rho_i\,\Delta r_i$, donde $\Delta r_i$ es la diferencia de radio entre capas consecutivas.


La \textbf{energía potencial gravitacional} de una configuración esférica autogravitante se calcula a partir del trabajo necesario para ensamblar la estrella capa por capa:
\begin{equation}
   d\Omega = -\frac{G\, M(r)\, dM(r)}{r},
\end{equation}
y la energía potencial total es:
\begin{equation}
   \Omega = -\int_0^R \frac{G\, M(r)}{r}\, dM(r).
\end{equation}
La relación entre la energía potencial y la masa y radio totales define el parámetro adimensional $q$:
\begin{equation}
   -\Omega_\star = q\,\frac{G M_\star^2}{R_\star}, \qquad q > 0,
\end{equation}
que se despeja directamente como:
\begin{equation}
   q = \frac{|\Omega_\odot|\, R_\odot}{G\, M_\odot^2}.
\end{equation}


\subsection*{Fracción de presión de radiación --- Proposición 7}
\label{sec:teoria-2b}


Para un plasma completamente ionizado el \textbf{peso molecular medio} depende de las fracciones de masa del hidrógeno $X$, el helio $Y$ y los metales $Z = 1-X-Y$:
\begin{equation}
   \mu = \frac{1}{2X + \dfrac{3}{4}Y + \dfrac{1}{2}Z}.
\end{equation}


los valores son usualmente de acuerdo con \cite{grassitelliSubsonicStructureOptically2018} son X = 0.71 y Y = 0.277 de forma que que Z = 0.013 que daría un valor promedio de $\mu = 0.612$, además observando la función ésta podrá variar entre 2 (solo otros metales) y 1/2 (solo hidrógeno)


La \textbf{proposición 7} del modelo estándar de Eddington establece la cota superior de la fracción de presión de gas $\beta^*$ en el centro de una estrella en equilibrio hidrostático. Definiendo la constante de radiación $a = 4\sigma/c$, la constante adimensional $B$ es:
\begin{equation}
   B = \frac{\pi\, a\, G^3 M_\odot^2\, m_h^4}{18\, k_B^4},
\end{equation}
y la fracción de presión de radiación central máxima satisface:
\begin{equation}
   \left(\frac{M_\star}{M_\odot}\right)^{\!2} B\,\mu^4 = \frac{1 - \beta^*}{\beta^{*4}}.
\end{equation}
El cambio de variable $x = 1 - \beta^*$ transforma la ecuación en:
\begin{equation}
   \frac{x}{(1-x)^4} = K, \qquad K = \left(\frac{M_\star}{M_\odot}\right)^{\!2} B\,\mu^4,
\end{equation}
función que es monótona creciente en $x\in(0,1)$ y tiene por tanto solución única para cada $K>0$, la cual se resuelve numéricamente hallando la raíz con el método de Brent (\texttt{scipy.optimize.brentq}).


El corte $1-\beta^*=0{,}5$ corresponde al caso en que la presión de radiación iguala a la del gas ($P_\mathrm{rad} = P_\mathrm{gas}$). Por encima de este límite, la configuración se aproxima al límite de Eddington: la proposición 7 asume que la presión de radiación es despreciable frente a la del gas para derivar el equilibrio hidrostático, supuesto que se rompe al llegar a $\beta^*=0{,}5$. Adicionalmente, con $1-\beta_c > 0{,}5$ en el centro, la fracción de presión de radiación aumenta hacia las capas externas, lo que puede inducir una pérdida violenta de masa o una inversión de densidad que destruye el equilibrio hidrostático de la estrella \cite{grassitelliSubsonicStructureOptically2018}; de ahí que $\beta^*=0{,}5$ actúe como cota máxima físicamente permitida en el interior estelar.


% ============================================================
\section*{Actividad 2A: Distribución de masa y energía potencial gravitacional}


Se cargaron los datos del Modelo S \cite{ChristensenDalsgaard1996} con 2482 capas radiales. se calcularon las diferencias de radio entre capas consecutivas, y la masa encerrada $M(r)$ se calculó como la suma acumulada de $\Delta M_i = 4\pi r_i^2\,\rho_i\,\Delta r_i$ (ecuaciones 1 y 2), usando la densidad en unidades SI obtenida tras la conversión de las unidades originales del modelo (CGS).


La figura \ref{fig:masa} muestra la masa encerrada $M(r)$ en función del radio relativo $r/R_\odot$. Se observa que la masa se concentra en la parte central ($r/R_\odot < 0{,}5$), lo que refleja el perfil de densidad fuertemente centralizado del Modelo S que corresponde con la teoría.


\begin{figure}[H]
   \centering
   \safeincludegraphics{mass_distribution.png}
   \caption{Masa encerrada $M(r)$ en función del radio relativo $r/R_\odot$ calculada a partir del Modelo S de Christensen-Dalsgaard et al.\ \cite{ChristensenDalsgaard1996}. Las constantes físicas empleadas provienen de CODATA 2014 vía \texttt{astropy} \cite{mohrCODATARecommendedValues2016}.}
   \label{fig:masa}
\end{figure}


La energía potencial gravitacional se calculó mediante la suma acumulada de $d\Omega_i = -G\,M(r_i)\,\Delta M_i / r_i$ (ecuaciones 3 y 4). El valor absoluto obtenido es:
\begin{equation*}
   |\Omega_\odot| = 6{,}12 \times 10^{41}\ \mathrm{J}.
\end{equation*}
Aplicando la ecuación (6), el parámetro $q$ resulta:
\begin{equation*}
   q = 1{,}617.
\end{equation*}
Este resultado es coherente con el valor esperado para una estrella politrópica de índice $n=3$ ($q\approx 1{,}5$); la desviación refleja que el perfil de densidad real del Sol no es exactamente politrópico.


% ============================================================
\section*{Actividad 2B: Cota superior de la fracción de presión de radiación}


Para la composición solar estándar ($X=0{,}71$, $Y=0{,}277$, $Z=0{,}013$) la ecuación (7) da $\mu_\odot = 0{,}612$. Dado que $\mu$ puede variar entre $0{,}5$ (hidrógeno puro totalmente ionizado) y $2{,}0$ (metales pesados totalmente ionizados), se exploraron siete valores: $\mu \in \{0{,}5;\, 0{,}612;\, 0{,}7;\, 1{,}0;\, 1{,}2;\, 1{,}5;\, 2{,}0\}$.


La constante $B$ se evaluó numéricamente usando los valores CODATA 2014 de \texttt{astropy} \cite{mohrCODATARecommendedValues2016}, obteniéndose $B \approx 3{,}35\times10^{-2}$. Para cada par $(M_\star/M_\odot,\,\mu)$ se resolvió la ecuación (10) con \texttt{scipy.optimize.brentq}, barriendo 500 valores de masa en el rango $[0{,}1,\;60]\,M_\odot$.


La figura \ref{fig:beta} muestra la fracción de presión de radiación $1-\beta^*$ en función de la masa estelar para cada valor de $\mu$. Se observa que la fracción de presión de radiación aumenta con la masa y con el peso molecular medio: composiciones más pesadas (mayor $\mu$, mayor metalicidad) producen estrellas intrínsecamente más densas, lo que permite sostener una fracción de presión de radiación más elevada. La línea discontinua en $1-\beta^*=0{,}5$ marca la cota máxima: por encima de ésta la presión de radiación domina sobre la del gas y la estrella, por lo que se rompe la suposición de equilibrio hidrostático, ya que depende de poder considerar la presión de radiación despreciable invalidando el modelo. Asimismo, se verifica que para masas superiores a $\approx 60\,M_\odot$ todas las configuraciones superan dicha cota independientemente del valor de $\mu$.


\begin{figure}[H]
   \centering
   \safeincludegraphics{radiation_pressure.png}
   \caption{Cota superior de la fracción de presión de radiación central $1-\beta^*$ en función de $M_\star/M_\odot$ para distintos valores del peso molecular medio $\mu$ (plasma totalmente ionizado). La línea discontinua indica el límite $\beta^*=0{,}5$ donde la presión de radiación iguala a la del gas, definiendo la cota máxima físicamente permitida en el interior estelar.}
   \label{fig:beta}
\end{figure}


% ============================================================
\bibliographystyle{plain}
\bibliography{biblography}


\end{document}

